\section{Scaling Laws Have a Mechanistic Interpretation}
\label{sec:scaling-laws}

Kaplan et al.\ and subsequent work established that loss scales predictably
as a power law with compute, parameters, and data~\cite{kaplan2020scaling}.
This is an empirical regularity. No principled derivation of the exponents
has been established, though several theoretical programs have attempted
partial accounts.

\subsection*{The BP Interpretation}

Under the BP interpretation, scaling adds inference capacity in well-defined
ways: more layers add rounds of BP, more $D_{ff}$ adds OR gate capacity,
more $D_{\mathrm{model}}$ adds superposition capacity, and more $W$ adds
factor graph nodes. Each of these is a smooth, monotone increase in a
well-defined quantity. Power-law scaling in loss is qualitatively consistent
with a system where each additional unit of capacity provides diminishing
but consistent marginal improvement in inference quality over a fixed
distribution.

\subsection*{The AND Head Invariant}

Empirically, boolean-AND head count stays roughly fixed within a model
family while hub and OR capacity scales. This is consistent with the BP
interpretation: the routing primitive is compact and does not need to scale,
while inference capacity does. This evidence is preliminary --- three model
families --- and should not be overstated.

\subsection*{Epistemic Status}

\begin{center}
\begin{tabular}{ll}
\toprule
Claim & Status \\
\midrule
Scaling dimensions have BP interpretations & Mechanistic interpretation \\
More layers enable deeper reasoning chains & Mechanistic interpretation \\
AND head invariance is consistent with BP & Empirical observation (preliminary) \\
BP explains power-law exponents & Not established \\
\bottomrule
\end{tabular}
\end{center}